\begin{abstract}
Neural image codecs typically decode an entire image with a fixed network depth, although the benefit of additional decoding layers can vary substantially across image regions. This paper introduces a spatially adaptive coding framework designed to exploit this variation. Instead of assigning the same synthesis complexity everywhere, our approach allocates decoder depth according to the reconstruction needs of individual regions. We investigate this idea in two forms. The first introduces spatial early exits into a shared decoder, where regions stop at different synthesis depths while reusing a common coded representation. The second selects among independently trained codecs with different decoder depths, allowing both the representation and synthesis network to adapt to each region. In this way, additional computation is spent only on regions where deeper decoding is useful. Building on these designs, we present \textbf{DepthLIC}, a learned image compression framework that turns decoder depth from a fixed architectural choice into a spatially adaptive coding decision.
\end{abstract}